% ── Rozwiązanie: Matrix-Chain-Order -- instancja 1 ────────────────
\subsection{Matrix-Chain-Order -- instancja 1}

Dane: $n=4$, $p=(5,\ 4,\ 6,\ 2,\ 7)$, czyli
\[
	A_1\ (5\times 4),\quad A_2\ (4\times 6),\quad A_3\ (6\times 2),\quad A_4\ (2\times 7).
\]
Wypełniamy tablice $m$ i~$s$ \emph{przekątnymi} po~rosnącej długości
podciągu $L$; w~każdym kroku dopisane komórki wyróżniono na
\textcolor{accentB}{pomarańczowo}, a~$\times$ to pola nieużywane
($i>j$ dla~$m$, $i\geq j$ dla~$s$). Skrót: $A_i=p_{i-1}\times p_i$.

\medskip
\noindent Stan początkowy ($L=1$, pojedyncze macierze, koszt $0$):\\[0.7em]
\begin{minipage}{0.45\linewidth}
	$m$
	$\begin{bNiceArray}{c|cccc}[margin,hvlines]
			i \backslash j & 1 & 2 & 3 & 4 \\
			1 & 0 &  &  &  \\
			2 & \times & 0 &  &  \\
			3 & \times & \times & 0 &  \\
			4 & \times & \times & \times & 0
		\end{bNiceArray}$
\end{minipage}\hfill
\begin{minipage}{0.45\linewidth}
	$s$
	$\begin{bNiceArray}{c|cccc}[margin,hvlines]
			i \backslash j & 1 & 2 & 3 & 4 \\
			1 & \times &  &  &  \\
			2 & \times & \times &  &  \\
			3 & \times & \times & \times &  \\
			4 & \times & \times & \times & \times
		\end{bNiceArray}$
\end{minipage}

\medskip
\noindent{\color{placeholdercolor}\rule{\linewidth}{0.4pt}}\par\smallskip
\noindent{\color{dygresjagray}\footnotesize$\begin{array}{c|*{5}{c}} i & 0 & 1 & 2 & 3 & 4 \\ \hline p_i & 5 & 4 & 6 & 2 & 7 \end{array}$\qquad $A_i = p_{i-1}\times p_i$}
\begin{flalign*}
	L=2: \quad & m(1, 2) = 5 \cdot 4 \cdot 6 = 120 &  & \\
	           & m(2, 3) = 4 \cdot 6 \cdot 2 = 48  &  & \\
	           & m(3, 4) = 6 \cdot 2 \cdot 7 = 84  &  &
\end{flalign*}

\medskip
\noindent po $L=2$ (pary sąsiednich macierzy):\\[0.7em]
\begin{minipage}{0.45\linewidth}
	$m$
	$\begin{bNiceArray}{c|cccc}[margin,hvlines]
			i \backslash j & 1 & 2 & 3 & 4 \\
			1 & 0 & \color{accentB}120 &  &  \\
			2 & \times & 0 & \color{accentB}48 &  \\
			3 & \times & \times & 0 & \color{accentB}84 \\
			4 & \times & \times & \times & 0
		\end{bNiceArray}$
\end{minipage}\hfill
\begin{minipage}{0.45\linewidth}
	$s$
	$\begin{bNiceArray}{c|cccc}[margin,hvlines]
			i \backslash j & 1 & 2 & 3 & 4 \\
			1 & \times & \color{accentB}1 &  &  \\
			2 & \times & \times & \color{accentB}2 &  \\
			3 & \times & \times & \times & \color{accentB}3 \\
			4 & \times & \times & \times & \times
		\end{bNiceArray}$
\end{minipage}

\medskip
\noindent{\color{placeholdercolor}\rule{\linewidth}{0.4pt}}\par\smallskip
\noindent{\color{dygresjagray}\footnotesize$\begin{array}{c|*{5}{c}} i & 0 & 1 & 2 & 3 & 4 \\ \hline p_i & 5 & 4 & 6 & 2 & 7 \end{array}$\qquad $A_i = p_{i-1}\times p_i$}
\begin{flalign*}
	L=3: \quad & m(1, 3) = \min \begin{cases}
		                            m(1, 1) + m(2, 3) + 5 \cdot 4 \cdot 2 \\
		                            m(1, 2) + m(3, 3) + 5 \cdot 6 \cdot 2
	                            \end{cases} = \min \begin{cases}
		                                               0 + 48 + 40 \\
		                                               120 + 0 + 60
	                                               \end{cases} = \min \begin{cases} 88 \\ 180 \end{cases} = 88 \\
	           & m(2, 4) = \min \begin{cases}
		                            m(2, 2) + m(3, 4) + 4 \cdot 6 \cdot 7 \\
		                            m(2, 3) + m(4, 4) + 4 \cdot 2 \cdot 7
	                            \end{cases} = \min \begin{cases}
		                                               0 + 84 + 168 \\
		                                               48 + 0 + 56
	                                               \end{cases} = \min \begin{cases} 252 \\ 104 \end{cases} = 104
\end{flalign*}

\medskip
\noindent po $L=3$:\\[0.7em]
\begin{minipage}{0.45\linewidth}
	$m$
	$\begin{bNiceArray}{c|cccc}[margin,hvlines]
			i \backslash j & 1 & 2 & 3 & 4 \\
			1 & 0 & 120 & \color{accentB}88 &  \\
			2 & \times & 0 & 48 & \color{accentB}104 \\
			3 & \times & \times & 0 & 84 \\
			4 & \times & \times & \times & 0
		\end{bNiceArray}$
\end{minipage}\hfill
\begin{minipage}{0.45\linewidth}
	$s$
	$\begin{bNiceArray}{c|cccc}[margin,hvlines]
			i \backslash j & 1 & 2 & 3 & 4 \\
			1 & \times & 1 & \color{accentB}1 &  \\
			2 & \times & \times & 2 & \color{accentB}3 \\
			3 & \times & \times & \times & 3 \\
			4 & \times & \times & \times & \times
		\end{bNiceArray}$
\end{minipage}

\medskip
\noindent{\color{placeholdercolor}\rule{\linewidth}{0.4pt}}\par\smallskip
\noindent{\color{dygresjagray}\footnotesize$\begin{array}{c|*{5}{c}} i & 0 & 1 & 2 & 3 & 4 \\ \hline p_i & 5 & 4 & 6 & 2 & 7 \end{array}$\qquad $A_i = p_{i-1}\times p_i$}
\begin{flalign*}
	L=4: \quad & m(1, 4) = \min \begin{cases}
		                            m(1, 1) + m(2, 4) + 5 \cdot 4 \cdot 7 \\
		                            m(1, 2) + m(3, 4) + 5 \cdot 6 \cdot 7 \\
		                            m(1, 3) + m(4, 4) + 5 \cdot 2 \cdot 7
	                            \end{cases} = \min \begin{cases}
		                                               0 + 104 + 140 \\
		                                               120 + 84 + 210 \\
		                                               88 + 0 + 70
	                                               \end{cases} = \min \begin{cases} 244 \\ 414 \\ 158 \end{cases} = 158
\end{flalign*}

\medskip
\noindent tablice końcowe:\\[0.7em]
\begin{minipage}{0.45\linewidth}
	$m$
	$\begin{bNiceArray}{c|cccc}[margin,hvlines]
			i \backslash j & 1 & 2 & 3 & 4 \\
			1 & 0 & 120 & 88 & \color{accentB}158 \\
			2 & \times & 0 & 48 & 104 \\
			3 & \times & \times & 0 & 84 \\
			4 & \times & \times & \times & 0
		\end{bNiceArray}$
\end{minipage}\hfill
\begin{minipage}{0.45\linewidth}
	$s$
	$\begin{bNiceArray}{c|cccc}[margin,hvlines]
			i \backslash j & 1 & 2 & 3 & 4 \\
			1 & \times & 1 & 1 & \color{accentB}3 \\
			2 & \times & \times & 2 & 3 \\
			3 & \times & \times & \times & 3 \\
			4 & \times & \times & \times & \times
		\end{bNiceArray}$
\end{minipage}

\paragraph{Optymalne nawiasowanie.}
Odtwarzamy je z~tablicy $s$, rekurencyjnie dzieląc $A_i\dots A_j$ w~miejscu
$s(i,j)$:
\[
	s(1,4)=3 \;\Rightarrow\; (A_1 A_2 A_3)(A_4),\qquad
	s(1,3)=1 \;\Rightarrow\; (A_1)(A_2 A_3).
\]
Stąd optymalne nawiasowanie
\[
	\boxed{\big(A_1\,(A_2\,A_3)\big)\,A_4}
\]
o~minimalnym koszcie $m(1,4)=158$ mnożeń skalarnych.
